\chapter{Tic-Tac-Toe without Revealing the Strategy}
\label{chap:tic-tac-toe-private-strategy}

\begin{quote}
\textbf{Author's Intuition.}

\emph{A perfect strategy in Tic-Tak-Toe is an undefeated one, one that might
win or tie any game. Here is a way to prove to your adversary that you know the
perfect strategy without revealing the strategy. Commit to a session seed;
avoid the adversary cloning the strategy; let the adversary select the initial
conditions; inspect all the Monte Carlo scenarios; investigate whether one can
invert the matrix; and always include code for demonstration.}
\end{quote}

The original intuition uses \emph{strategy} broadly and uses \emph{perfect} to
mean undefeated. The formal object below is a deterministic positional
\emph{policy}. The property checked is \emph{non-loss from one fixed initial
state}, not local minimax optimality at every state. This translation preserves
the question while preventing the words from carrying a stronger theorem than
the construction.

These notes contain several distinct research questions. This chapter
preserves them, but it does not present them as one completed mechanism. A
session identifier, commitment randomness, and a policy-generation seed are
different objects. Random sampling is not an exhaustive proof. No matrix has
yet been defined whose inverse implies anything about Tic-Tac-Toe. These are
open directions, not hidden steps in the construction below.

\begin{center}
\fbox{\begin{minipage}{0.91\textwidth}
\textbf{Current Result, in Plain Words.}

\begin{enumerate}
    \item We can disclose one complete policy and exhaustively check that it
    never loses under the fixed game model.
    \item We can commit to the exact policy bytes and later verify a disclosed
    opening under a fixed educational profile.
    \item We cannot yet prove non-loss while keeping that policy secret. No
    current circuit, commitment, or test is a zero-knowledge proof.
\end{enumerate}
\end{minipage}}
\end{center}

\begin{abstract}
This chapter separates a private-strategy aspiration from the public predicate
that any later private proof would have to establish. For ordinary
three-by-three Tic-Tac-Toe, we define legal reachable states, a canonical
deterministic positional policy, and a non-loss condition for a fixed player
role. A plain C11 diagnostic checker validates a fully disclosed policy and
exhaustively branches over every legal opponent move from the standard empty
board. A second C11 evaluator expresses the same verdict as a fixed bottom-up
recurrence over the reachable game DAG, without recursion or a
policy-selected successor address. A third dependency-free C11 component emits
and evaluates fixed XOR/AND circuits for the game predicate
\(\mathsf{Core}\). A fourth component composes the existing educational
BLAKE2b-256 commitment with an exact Tic-Tac-Toe context, checks disclosed
openings, and counts one frozen XOR/AND representation of
\(\mathsf{Bind}\) without materializing it. All current paths remain
disclosure baselines, not a solution to
the chapter title. We distinguish
knowledge of a committed policy, policy legality, global non-loss, private
per-state evaluation, and resistance to extraction through repeated queries.
We then state a candidate committed-policy relation for future zero-knowledge
argument-of-knowledge research, compare bounded implementation routes, and
complete the backend-neutral Boolean-circuit step for
\(\mathsf{Core}\), freeze an external policy-commitment profile, and measure
the direct Boolean cost of \(\mathsf{Bind}\). We prove why the same policy
witness must satisfy both predicates and preserve the measured direct profile
as a comparison baseline. We do not emit that large circuit or select an
unspecified ``proof-friendly'' replacement before choosing the proof
representation that would consume it. Every proof-internal commitment and
every proof transformation remain open. We identify the leakage, abort,
composition, and setup questions that remain open. The resulting
Core--Bind--proof separation is the book's canonical running example; its
assumptions and conclusions do not transfer to another problem merely because
that problem cites it.
\end{abstract}

\section{Cooperation Problem}
\label{sec:ttt-cooperation-problem}

A player may wish to demonstrate that it possesses a strategy that cannot be
defeated while retaining the particular strategy as private information. The
verifier wants evidence stronger than one successful game, but the player does
not want to disclose a complete response table.

The parties are only partially aligned. Both benefit from a precise and
checkable claim. The claimant benefits from keeping its policy private, while
the verifier benefits from finding a losing branch or extracting useful
responses. The narrow research question is:

\begin{quote}
Can a claimant prove knowledge of one fixed, legal, non-losing policy for a
specified role and initial state without revealing that policy beyond the
truth of the claim?
\end{quote}

The first public slice answered a prerequisite question instead: what exact
predicate should be checked if the policy is disclosed? The second fixed its
game-DAG recurrence; later slices froze that recurrence as public Boolean
netlists, fixed the external commitment bytes, and measured the direct cost of
recomputing that commitment inside the relation. The present decision keeps
that direct profile as a baseline but declines to materialize it before a
proof representation is selected. None of these steps hides the policy.
Defining the predicate and binding contract first prevents a future proof
system from hiding an ambiguous or incorrect computation.

\section{Disclosure and Trusted-Evaluator Baselines}
\label{sec:ttt-baselines}

Under full disclosure, the claimant sends the complete strategy to the
verifier. The verifier can check every entry and simulate every opponent
continuation. This provides transparent verification but no policy privacy.

A trusted evaluator offers another baseline. The claimant privately supplies
the policy; the evaluator checks it and announces a verdict. This hides the
policy from the verifier only by moving trust to the evaluator, who learns the
policy and may misreport, retain, or redistribute it.

The public checker in this chapter implements the first baseline. A future
cryptographic construction would seek some of the evaluator's confidentiality
without inheriting its discretionary power. Cryptography is relevant only
after the statement, witness, verifier, and leakage have been fixed.

\section{Why This Is the Book's Running Example}
\label{sec:ttt-running-example}

Tic-Tac-Toe is useful here precisely because it is not mysterious. Its rules
are public, its state space is finite, a losing branch can be replayed, and a
reader can reproduce the complete semantic check. The strategy is not a scarce
real-world secret: anyone can construct a non-losing policy. The example is a
laboratory for protocol anatomy, not a claim that Tic-Tac-Toe knowledge has
economic value.

Table~\ref{tab:ttt-concept-ladder} is the reference map used throughout this
chapter and intended for later chapters. Each layer answers a different
question. Skipping a layer does not cause the next layer to supply it.

\begin{table}[H]
\centering
\footnotesize
\begin{tabular}{@{}
>{\raggedright\arraybackslash}p{0.16\textwidth}
>{\raggedright\arraybackslash}p{0.19\textwidth}
>{\raggedright\arraybackslash}p{0.27\textwidth}
>{\raggedright\arraybackslash}p{0.27\textwidth}@{}}
\hline
Layer & Meaning here & What it establishes & What it does not establish \\
\hline
Game model & Rules and reachable states & Meaning of legal play & A policy \\
Policy $\pi$ & Complete deterministic response table & Which move is selected & Legality, non-loss, or privacy \\
$\mathsf{Core}$ & Canonical legality and global non-loss predicate & The semantic property of $\pi$ & A connection to commitment $C$ \\
Commitment $C$ & Context-bound digest of policy bytes & Intended byte binding under the stated assumptions & Policy quality, knowledge, or chronology \\
$\mathsf{Bind}$ & Opening relation for $C$ & The witness opens $C$ in the stated context & Legality or non-loss \\
$R=\mathsf{Core}\land\mathsf{Bind}$ & The same policy satisfies both halves & The exact committed-policy statement & A proof protocol \\
Argument of knowledge & Future protocol for the selected relation & Verifier acceptance plus a stated extraction property & Witness privacy by itself \\
Zero knowledge & Privacy property of that protocol & A stated bound on transcript leakage & Identity, fairness, or delivery \\
Lifecycle & Publication, reuse, opening, abort, and expiration & Real execution meaning & Supplied by none of the layers above \\
\hline
\end{tabular}
\caption{The Tic-Tac-Toe concept ladder.}
\label{tab:ttt-concept-ladder}
\end{table}

Later chapters may reuse this ladder only after replacing every
problem-specific part: parties, canonical hidden object, public model,
predicate, leakage, adversary, setup, and lifecycle. They may cite the
shared-witness lesson in Section~\ref{sec:ttt-shared-witness} and the oracle
leakage result in Section~\ref{sec:ttt-repeated-query}; they inherit no
Tic-Tac-Toe security conclusion.

\section{Fixed Game Model}
\label{sec:ttt-model}

The game is ordinary three-by-three Tic-Tac-Toe. Number the cells from $0$
through $8$ in row-major order. Each cell is empty, X, or O. X moves first;
the roles alternate; a row, column, or diagonal of one mark wins; and play
stops immediately after a win. A full board without a winner is a draw. The
initial state $s_0$ is the empty board.

A state is \emph{legal and reachable} only if it occurs along such a play from
$s_0$. In particular, boards with impossible mark counts, simultaneous
winners, or moves after a win are not reachable game states.

Encode cell values as

\[
    \mathsf{EMPTY}=0,\qquad \mathsf{X}=1,\qquad \mathsf{O}=2,
\]

and encode a board $b$ by

\[
    \operatorname{index}(b)=\sum_{i=0}^{8} b_i 3^i.
\]

For orientation, the cell numbering and one reachable board are

\begin{center}
\begin{tabular}{c|c|c}
0 & 1 & 2 \\ \hline
3 & 4 & 5 \\ \hline
6 & 7 & 8
\end{tabular}
\qquad
\begin{tabular}{c|c|c}
X & $\cdot$ & $\cdot$ \\ \hline
$\cdot$ & O & $\cdot$ \\ \hline
$\cdot$ & $\cdot$ & $\cdot$
\end{tabular}
\end{center}

The board on the right has index
$1\cdot3^0+2\cdot3^4=163$. It is X's turn, so a policy for claimant X
stores its selected next cell at byte $\pi[163]$. This one lookup is the small
unit repeated across the complete table.

Thus there are $3^9=19{,}683$ possible byte-addressable board encodings,
although most do not describe legal reachable states. Exhaustive enumeration
from the empty board, without expanding terminal states, gives 5,478 reachable
boards: 2,423 nonterminal X-to-move boards, 2,097 nonterminal O-to-move boards,
and 958 terminal boards. The C tests use these independently expected counts
as finite-model cross-checks, not as cryptographic evidence.

\section{Canonical Deterministic Policy}
\label{sec:ttt-policy}
\label{def:ttt-policy-v1}

Fix a claimant role $p\in\{\mathsf{X},\mathsf{O}\}$. Let $S_p$ be the set
of legal reachable nonterminal states in which $p$ moves, and let $A(s)$
be the empty cells in state $s$. A deterministic positional policy is a
function

\[
    \pi:S_p\longrightarrow\{0,\ldots,8\},
    \qquad \pi(s)\in A(s)\text{ for every }s\in S_p.
\]

Canonical Policy Version~1 represents this function as a 19,683-byte table.
At an index corresponding to a state in $S_p$, the byte is one legal cell
number from $0$ to $8$. Every invalid board, opponent-turn board, terminal
board, and other non-required position contains the distinguished byte
\texttt{0xff}. Consequently, the encoding has no semantically ignored bytes.

The claimant role is a separate public input. The table does not encode the
claimant's identity, authorship, originality, randomness, session, or promise
to follow the policy later.

The original note proposes committing to a session seed. A seed alone commits
to nothing. A binding commitment to that seed can bind the selected policy
only if a public deterministic algorithm

\[
    \pi=\mathsf{GeneratePolicy}(\mathit{seed})
\]

is fixed and that equality is included in the checked relation. Without such
a rule, a seed and a policy are unrelated strings. A unique public session
identifier distinguishes execution context only when it is incorporated into
the committed message, public statement, or authenticated transcript;
commitment randomness is private opening material used by a hiding
construction; and a generation seed may compactly select a policy. None
substitutes for the others.

\section{Legal, Non-Losing, and Minimax-Optimal}
\label{sec:ttt-policy-properties}

These terms must not be merged.

\begin{itemize}
    \item A \emph{legal} policy selects an empty cell at every required
    claimant-turn state and has the canonical no-move byte everywhere else.
    \item A policy is \emph{non-losing from $s_0$} if, when it is followed
    at every claimant turn and the opponent may make any legal move, every
    terminal branch is a claimant win or a draw.
    \item A policy is \emph{locally minimax-optimal} if its choice has optimal
    game-theoretic value at every state under consideration.
\end{itemize}

For ordinary Tic-Tac-Toe, the value of the empty board under optimal play is a
draw. A non-losing policy therefore realizes the optimal value from $s_0$,
but it need not choose an optimal move at every legal off-path state. The
current checker establishes non-loss only from the fixed empty board. It does
not establish local optimality at every state, uniqueness, or superiority to
another non-losing policy.

Allowing a verifier to choose arbitrary initial conditions changes the claim.
Some legal states are already forced losses for one role. A challenged-state
protocol must define a public allowed set, include the initial state in the
statement, and distinguish a bad continuation from a position whose value was
already losing. This remains open.

\section{Public Exhaustive Predicate}
\label{sec:ttt-public-predicate}
\label{def:ttt-public-nonlose}

For a role $p$ and canonical policy $\pi$, define
$\mathsf{PublicNonLose}(p,\pi)=1$ exactly when:

\begin{enumerate}
    \item every required entry is present and selects an empty cell;
    \item every non-required entry is \texttt{0xff}; and
    \item from $s_0$, following $\pi$ at claimant turns and branching over
    every legal opponent move never reaches an opponent win.
\end{enumerate}

The verifier first classifies the complete legal reachable-state graph. It
rejects missing, out-of-range, occupied-cell, and unexpected entries. It then
performs a deterministic depth-first traversal. At a claimant turn it follows
the policy's single move. At an opponent turn it explores every empty cell in
ascending order. It stops at every win or full-board draw. On the first
opponent win it returns the complete chronological move sequence.

This is exhaustive enumeration, not Monte Carlo sampling. A random sample can
miss the sole losing branch of an otherwise successful policy, and therefore
cannot by itself establish the universal no-loss predicate. Sampling could
support a separately stated statistical test under a named distribution and
error bound, but that would be a different claim.

The original matrix-inversion note is also not part of this predicate. No
matrix, encoding, or theorem currently connects invertibility to legal play or
non-loss. If such a connection is discovered, it should first be stated as a
separate mathematical relation and compared with direct exhaustive checking.

\section{C Implementation Companion and Evidence}
\label{sec:ttt-implementation}

The executable companion keeps four artifacts separate so their evidence
cannot be confused:

\begin{table}[H]
\centering
\small
\begin{tabular}{@{}
>{\raggedright\arraybackslash}p{0.19\textwidth}
>{\raggedright\arraybackslash}p{0.31\textwidth}
>{\raggedright\arraybackslash}p{0.39\textwidth}@{}}
\hline
Artifact & Teaching purpose & Boundary \\
\hline
Recursive checker & Explains rejection with a replayable losing line & Discloses the complete policy and uses diagnostic recursion \\
Fixed DAG & Evaluates the same predicate on a public schedule & Still a plaintext C evaluation, not a circuit proof \\
Core circuits & Encode canonical validity and non-loss as XOR/AND gates & Witness-typed inputs remain plaintext during evaluation \\
Bind companion & Checks a disclosed commitment opening and measures one direct circuit shape & Does not check policy quality or produce a private proof \\
\hline
\end{tabular}
\caption{The four current executable artifacts.}
\label{tab:ttt-implementation-artifacts}
\end{table}

\begin{sloppypar}
The preserved public-checker interface is \path{include/ac/ttt.h}; both
evaluators remain in \path{src/protocols/ttt.c}; and their tests remain in
\path{tests/test_ttt.c}. The Boolean interface is split between
\path{include/ac/bool_circuit.h} and
\path{include/ac/ttt_core_circuit.h}; its implementations are
\path{src/circuits/bool_circuit.c} and
\path{src/circuits/ttt_core_circuit.c}; and its tests are
\path{tests/test_bool_circuit.c} and
\path{tests/test_ttt_core_circuit.c}. The demonstration is
\path{src/tic_tac_toe_without_revealing_the_strategy/ttt_demo.c}. The protected
checker interface, implementation, and original 13-test suite remain
byte-identical independent baselines. These three public-predicate paths use C11,
caller-owned storage, and the standard library only. They perform no heap
allocation, randomness, hashing, commitment, network communication, or
libsodium call.
\end{sloppypar}

\begin{sloppypar}
The separate external-commitment interface is
\path{include/ac/ttt_bind.h}; its wrapper and symbolic measurement are in
\path{src/protocols/ttt_bind.c}; its focused tests are
\path{tests/test_ttt_bind.c}; and its warning-bearing driver is
\path{src/tic_tac_toe_without_revealing_the_strategy/ttt_bind_demo.c}.
This component reuses the repository's existing
\path{src/primitives/commitment.c} and therefore links libsodium for
BLAKE2b-256, random nonce generation, comparison, and clearing. It adds no
cryptographic library. The protected dependency-free \(\mathsf{Core}\)
targets remain separately linkable without libsodium.
\end{sloppypar}

Two deterministic fixtures serve different lessons. The original synthetic
Bind vector intentionally opens to a policy that fails Core, demonstrating
that opening consistency does not imply strategy quality. The positive vector
in \path{test-vectors/tic-tac-toe-reference-v1.txt} uses the valid X reference
policy. One grouped test carries those exact bytes through the recursive
checker, fixed DAG, actual Core-X Boolean evaluation, and Bind opening. The
Bind demo also has a deterministic reference mode, invoked with
\texttt{--reference-vector}. Its default mode still obtains a fresh session
and nonce. The fixed nonce is
public test material and supplies no hiding or deployment evidence.

Both evaluators distinguish API execution from the mathematical verdict. Null
pointers and invalid roles are API errors. A malformed but readable policy is
a successfully computed invalid-policy verdict. A canonical policy with an
opponent-winning continuation is a losing verdict accompanied by a replayable
trace from the recursive checker. The fixed-DAG report deliberately contains
no invented trace or DFS exploration count. The public reference-policy
generator uses deterministic backward
induction; it is a convenience for demonstrations and tests, not a secret or
cryptographic mechanism.

The test suite checks both roles, an independent graph and history census,
malformed encodings, late losing branches, deterministic repeated reports,
terminal-state handling, every legal one-entry alternative of both generated
reference policies, and a deterministic set of multi-entry policies. The
single-entry sweep contains 6,208 X policies and 5,439 O policies. It is
exhaustive only at Hamming distance one from those two reference tables, not
over all canonical policies. Such tests establish enumerated implementation
behaviour for tested builds. The induction argument below establishes the
mathematical recurrence; neither it nor the tests are machine-checked formal
verification of the C program. They do not establish policy privacy, zero
knowledge, knowledge extraction, authorship, or future behaviour.

Most importantly, the checker receives the entire 19,683-byte table. It does
not implement the title's non-revelation goal.

\section{Five Claims That Must Remain Separate}
\label{sec:ttt-five-claims}

\subsection{Knowledge of a Policy or of a Committed Policy}

The one-shot existential statement has policy witness $\pi$. The committed
statement has witness $(\pi,r)$, because the extractor must recover both the
policy and an opening accepted by the named relation. A commitment-shaped
value does not demonstrate knowledge of either witness. Formal proof of
knowledge requires an extractor-based definition relating verifier acceptance
to recovery of the stated witness \cite{bellare1993definingknowledge}. The
three public checker/Core paths have no commitment or extractor. The separate
external wrapper verifies disclosed openings but still has no extractor.

\subsection{Legality}

Once a complete policy is disclosed, legality is an ordinary deterministic
property. An unopened commitment cannot be scanned for missing, occupied, or
unexpected moves. Merely committing to bytes proves none of their semantics.

\subsection{Non-Loss}

Once a legal policy is disclosed, finite exhaustive traversal establishes the
implemented no-loss predicate. One successful game, a set of favourable
examples, or an unspecified Monte Carlo loop does not cover every adversarial
branch.

\subsection{Private Per-State Evaluation}

Returning $\pi(s)$ necessarily reveals that output. If the state and output
are public, a future claimant might return the move with a zero-knowledge proof
that it is consistent with the same committed policy. If the verifier's query
state must also be hidden, the problem becomes private two-party computation.
Neither mechanism is present here.

\subsection{Resistance to Cloning}

A reusable interface that returns the exact deterministic move on arbitrary
states exposes the table one entry at a time. Cryptography cannot conceal an
answer that the interface deliberately releases. Query restrictions,
statefulness, fresh session policies, or randomized policies may change how
quickly useful information accumulates, but each requires its own leakage
definition.

\section{Two Statement Modes and a Future Private Relation}
\label{sec:ttt-future-relation}

Commitment schemes formalize hiding and binding properties under stated
assumptions; Naor's construction is one foundational reference
\cite{naor1991bitcommitment}. That work does not prove the repository's
instructional hash commitment or the Tic-Tac-Toe predicate.

The chapter has two distinct statement modes. In \emph{existential mode}, no
external commitment is needed. With public statement
$x_{\exists}=(v_G,p,s_0)$ and witness $w_{\exists}=\pi$, define

\begin{equation}
\label{eq:ttt-existential-relation}
R_{\exists}(x_{\exists},w_{\exists})=
\mathsf{Core}(v_G,p,s_0,\pi).
\end{equation}

This relation defines which policy witnesses satisfy Core. A future argument
of knowledge for the relation could support the corresponding knowledge claim.
The relation alone does not do so, and it does not say that the policy was
fixed earlier or will be reused later.

In \emph{committed mode}, a future protocol may use the public instance

\[
 x=(v_G,p,s_0,C,\mathit{ctx})
\]

and the private witness

\[
 w=(\pi,r),
\]

where $v_G$ fixes the game and policy versions, $s_0$ fixes the initial state,
$C$ is a policy commitment, $\mathit{ctx}$ fixes the protocol, session, and
statement-round context, and $r$ is its opening information. Define the
relation

\begin{equation}
\label{eq:ttt-committed-policy-relation}
 \begin{aligned}
 R(x,w)=1 \quad\Longleftrightarrow\quad{}
 &v_G=1 \;\land\; p\in\{\mathsf{X},\mathsf{O}\}
     \;\land\; s_0=\mathsf{EmptyBoard} \\
 &{}\land\;\mathsf{Open}\bigl(C,\mathit{ctx},
     \mathsf{Encode}(v_G,p,s_0,\pi),r\bigr)=1 \\
 &{}\land\;\mathsf{PublicNonLose}(p,\pi)=1.
 \end{aligned}
\end{equation}

Here one policy $\pi$ must satisfy both semantic and opening conditions. The
symbol $v_G$ is shorthand for the fixed Version~1 game and policy semantics,
not a complete deployment version registry. A concrete private protocol must
separately identify its game rules, policy encoding, relation, external
commitment profile, proof system, and transcript versions. Adding those fields
to the Version~1 bytes would require a new profile; this chapter does not
silently reinterpret the existing digest.

\subsection{Commitment Publication Lifecycle}
\label{sec:ttt-commitment-lifecycle}

A digest has temporal meaning only within an externally enforced lifecycle.
Table~\ref{tab:ttt-commitment-lifecycle} states the minimum contract. No current
C API implements the ordered record, authentication, replay database, or
expiration policy.

\begin{table}[H]
\centering
\small
\begin{tabular}{@{}
>{\raggedright\arraybackslash}p{0.16\textwidth}
>{\raggedright\arraybackslash}p{0.33\textwidth}
>{\raggedright\arraybackslash}p{0.40\textwidth}@{}}
\hline
Stage & Required event & Meaning and boundary \\
\hline
Allocate & Fix versions, role, session, round, and statement purpose & Context bytes distinguish statements but do not authenticate them \\
Publish & Place $(C,\mathit{ctx})$ in an authenticated ordered record before the event it is meant to precede & Without this record, ``fixed earlier'' is not established \\
Prove & Bind the proof transcript to the exact recorded statement and use one shared witness & A proof for another digest, context, or policy is irrelevant \\
Reuse & Refer to the same recorded value under an explicit reuse policy & Reuse is not inferred from equal-looking local transcripts \\
Open or close & Verify an optional opening, or record withdrawal, abort, or expiration & Cryptography does not explain an abort or force an opening \\
\hline
\end{tabular}
\caption{External lifecycle required for a temporal commitment claim.}
\label{tab:ttt-commitment-lifecycle}
\end{table}

The current session and statement-round fields supply domain context. Only an
application or institution can supply the ordered publication event and decide
who may write, read, challenge, reuse, revoke, or close it.

Zero knowledge asks that the verifier learn no more from a proof than what is
already implied by the public statement. The foundational definition is due
to Goldwasser, Micali, and Rackoff
\cite{goldwasser1989knowledgecomplexity}. Goldreich, Micali, and Wigderson show
the conceptual breadth of zero-knowledge proofs for NP under their stated
assumptions \cite{goldreich1991proofsvalidity}. This existence result does not
select a practical backend, prove knowledge automatically, or establish the
security of a concrete composition.

\subsection{Target Security Experiments (Not Achieved)}
\label{sec:ttt-security-target}

The eventual object must be a family indexed by a security parameter
$\lambda$, with an explicit setup algorithm
$\mathsf{Setup}(1^\lambda)$ and a versioned relation family $R_\lambda$.
For the selected existential or committed mode, the protocol must state the
following; PPT denotes probabilistic polynomial-time.

\noindent\textbf{Completeness.}
For every $\lambda$, every public-parameter string
$\mathit{pp}\leftarrow\mathsf{Setup}(1^\lambda)$, and every $(x,w)$ satisfying
$R_\lambda(x,w)=1$, an honest prover using $\mathit{pp}$ makes an honest
verifier accept except with a stated completeness error. If the setup model
permits parameters not honestly sampled, their validation rule must be stated
separately.

\noindent\textbf{Argument of knowledge.}
For every PPT malicious prover and auxiliary input
that make the verifier accept statement $x$ above the selected knowledge error,
there must exist an efficient extractor under the selected extractor interface.
Except with a stated extraction-failure probability, it must output a witness
$w$ satisfying $R_\lambda(x,w)=1$: $\pi$ in existential mode or $(\pi,r)$ in
committed mode.

\noindent\textbf{Zero knowledge.}
For every PPT malicious verifier and auxiliary input $z$, there must exist a
PPT simulator given
$(1^\lambda,\mathit{pp},x,z)$ and the declared public outputs. Its simulated
view must be computationally indistinguishable from the verifier's real view,
under the selected setup and composition model.

These are target experiments, not theorems about current code. The fixed
32-byte Version~1 BLAKE2b profile is one concrete educational baseline and is
not presented as an asymptotic instantiation for arbitrary $\lambda$. The
author's full technical goal requires both the argument-of-knowledge and
zero-knowledge properties under named setup, adversary, commitment, channel,
and composition assumptions. This chapter implements none of that machinery.

\section{Selecting a Private-Proof Route}
\label{sec:ttt-proof-route}

The immediate research target is one global argument of knowledge for the
committed relation $R$, with the zero-knowledge target from
Section~\ref{sec:ttt-security-target}; it is not a service that returns policy
moves. The public instance and accept/reject outcome remain visible. A
per-state interface instead returns $\pi(s)$ and therefore reveals one policy
output on every query. Private per-state evaluation is a separate
two-party-computation problem with its own input privacy, output leakage,
fairness, and repeated-query requirements.

This target is provisional and concerns one standalone interactive execution
by classical probabilistic polynomial-time algorithms. For the combinatorial
soundness calculation below, the honest verifier samples independent uniform
challenges after the first proof-internal commitments. The malicious-verifier
zero-knowledge target must instead handle arbitrary polynomial-time challenge
strategies under the selected transformation. Human identity, transport
authentication, delivery, timing, concurrency, resets, adaptive corruption,
composition, endpoint compromise, fairness, and recovery remain outside this
initial target.

The relation also has two costs that must remain visible. For the direct
Version~1 baseline, write

\begin{equation}
\label{eq:ttt-v1-relation}
 R_{\mathrm{V1}}(x,w)=\mathsf{Core}(v_G,p,s_0,\pi)
        \land \mathsf{Bind}(C,\mathit{ctx},v_G,p,s_0,\pi,r),
\end{equation}

The predicate $\mathsf{Core}$ checks the game version, valid role, and empty
initial state. It also evaluates $\mathsf{PublicNonLose}$, which already
includes complete canonical-policy validation. The predicate $\mathsf{Bind}$
checks the exact commitment opening and context. Proving only $\mathsf{Core}$
does not connect the policy to $C$. Proving only $\mathsf{Bind}$ says nothing
about legality or non-loss.
Dropping $C$ entirely would merely prove knowledge of some non-losing ordinary
Tic-Tac-Toe policy, a capability that is already public and easy to reproduce.

\subsection{Why the Witness Must Be Shared}
\label{sec:ttt-shared-witness}

\begin{center}
\fbox{\begin{minipage}{0.88\textwidth}
\textbf{Reusable lesson.} A commitment must be consumed by the relation, and
the same witness must satisfy both halves.

\[
\begin{array}{ll}
\text{Wrong:} &
\mathsf{Core}(\pi_1)=1
\quad\text{and separately}\quad
\mathsf{Bind}(C,\pi_2,r)=1, \\[2mm]
\text{Required:} &
\mathsf{Core}(\pi)=1
\quad\land\quad
\mathsf{Bind}(C,\pi,r)=1.
\end{array}
\]
\end{minipage}}
\end{center}

The public digest cannot bind the proved policy to the commitment merely by
appearing in the statement. Let $y=(v_G,p,s_0)$ and define the Core-only lifted
relation

\[
 R_{\mathsf{Core}}((y,C,\mathit{ctx}),\pi)
   =\mathsf{Core}(y,\pi).
\]

\begin{proposition}[Unused commitment]
\label{prop:ttt-unused-commitment}
\mbox{}\par
Fix any policy $\pi$ and any values
$C,C',\mathit{ctx},\mathit{ctx}'$. Then

\[
 R_{\mathsf{Core}}((y,C,\mathit{ctx}),\pi)
 =R_{\mathsf{Core}}((y,C',\mathit{ctx}'),\pi).
\]

Thus acceptance by this relation cannot establish that $\pi$ opens $C$.
\end{proposition}

\textbf{Proof.}
By definition, both sides equal $\mathsf{Core}(y,\pi)$; neither $C$ nor its
context is read. $\square$

Consequently, even a perfectly sound argument of knowledge for this lifted
relation may correctly extract a non-losing policy while proving nothing about
$C$. That would be a relation-specification error, not a failure of the proof
system or the commitment.

The second failure is subtler. Define

\[
 \begin{aligned}
 G_y&=\{\pi:\mathsf{Core}(y,\pi)=1\},\\
 O_{C,\mathit{ctx},y}
   &=\{\pi:\exists r\;\mathsf{Bind}(C,\mathit{ctx},y,\pi,r)=1\}.
 \end{aligned}
\]

The committed-policy claim is
$G_y\cap O_{C,\mathit{ctx},y}\ne\varnothing$. A Core-only proof establishes
only $G_y\ne\varnothing$. Even separate proofs that $G_y$ and
$O_{C,\mathit{ctx},y}$ are nonempty do not suffice: two nonempty sets need not
intersect.

\begin{proposition}[Shared-witness linkage is necessary]
\label{prop:ttt-shared-witness}
Fix $y$, $C$, and $\mathit{ctx}$. Suppose there are policies $\pi_+$ and
$\pi_-$ and an opening $r_-$ such that

\[
 \mathsf{Core}(y,\pi_+)=1,\qquad
 \mathsf{Core}(y,\pi_-)=0,\qquad
 \mathsf{Bind}(C,\mathit{ctx},y,\pi_-,r_-)=1,
\]

and suppose $C$ has no accepting opening to any different policy byte string.
Then the Core and Bind existential statements are each true, but there is no
single $\pi,r$ satisfying their conjunction.
\end{proposition}

\textbf{Proof.}
The witness $\pi_+$ establishes the Core existential statement, while
$(\pi_-,r_-)$ establishes the Bind existential statement. The last assumption
forces every policy in $O_{C,\mathit{ctx},y}$ to equal $\pi_-$. Since
$\pi_-$ fails Core, $G_y\cap O_{C,\mathit{ctx},y}$ is empty. $\square$

The actual hash profile does not establish the proposition's last condition as
a theorem. Its published synthetic vector illustrates the mismatch: its
opening verifies while its policy fails Core, whereas a reference policy
passes Core. Given that known invalid opening, any efficiently produced
full-relation opening with a different, non-losing policy would constitute an
alternate-opening attack for the same digest and context. The profile's
informal computational binding discussion assumes such an attack is
infeasible; it supplies no formal reduction or numerical failure bound. The
test demonstrates the two code paths, not the assumption.

A future construction must share the exact policy, game-version, role, and
initial-state wires between both predicates, or enforce their equality
explicitly.

The external $C$ is not logically required for the one-shot claim ``I know
some non-losing policy.'' A proof of knowledge for Core would express that
weaker statement. A public commitment, or a formally specified native witness
commitment serving the same role, becomes necessary when the claim must refer
to one particular policy fixed earlier, reused across declared proofs, or
opened later. If $C$ is created alongside one proof and nothing else ever
refers to it, it contributes little temporal meaning while retaining its
relation cost. The current session and round fields separate contexts; they do
not themselves prove chronology or authenticate a person.

\subsection{Bounded Comparison}

The foundational generic result of Goldreich, Micali, and Wigderson shows why
an NP relation can be approached through zero knowledge
\cite{goldreich1991proofsvalidity}. A generic reduction remains valuable
theory, but implementing such a reduction as this chapter's companion would
add substantial representation and communication machinery without teaching
the structure of the finite-game predicate.

The MPC-in-the-head method of Ishai, Kushilevitz, Ostrovsky, and Sahai converts
a suitable multiparty computation into a zero-knowledge protocol
\cite{ishai2007zeroknowledgemp}. ZKBoo develops a concrete version for Boolean
circuits using symmetric primitives \cite{giacomelli2016zkboo}. Proposition
4.2 of ZKBoo analyzes its basic three-message protocol in the
commitment-hybrid model, assuming a correct, 2-private $(2,3)$-decomposition of
the circuit computation. In that model it is a Sigma protocol with special
honest-verifier zero knowledge and 3-special soundness. These properties do not
automatically provide malicious-verifier zero knowledge.

The proof-internal commitments to simulated views are distinct from the public
policy commitment $C$ checked by $\mathsf{Bind}$. A concrete experiment must
separately justify the binding and hiding of the view commitments, any
pseudorandom-tape assumptions, their encodings and domains, and their transcript
cost. Naming ZKBoo makes neither commitment layer compatible with the current
instructional hash commitment.

For one idealized repetition, a uniformly sampled challenge in
$\{1,2,3\}$ leaves the combinatorial soundness term at most $2/3$. With
independent repetitions and independently uniform challenges chosen after the
first messages, that term is $(2/3)^r$; the smallest integer making only this
term at most $2^{-128}$ is

\[
 r=\left\lceil\frac{128}{\log_2(3/2)}\right\rceil=219.
\]

This equation is not a total soundness bound or proof-size estimate for our
relation. A concrete argument must also include the binding-failure term of its
chosen view commitment. The Boolean cost of \(\mathsf{Core}\) is now measured
for the frozen construction below. One direct, deliberately unoptimized
Boolean representation of \(\mathsf{Bind}\) is also counted below, but not
emitted. Proof-internal commitments, simulated views, transformations, and
transcripts remain unselected or unmeasured. Neither a
\(\mathsf{Core}\) count nor a \(\mathsf{Core}\land\mathsf{Bind}\) count is a
proof-size estimate. Transformations beyond the basic Sigma protocol are required for the chapter's
malicious-verifier zero-knowledge and knowledge target. Making the protocol
non-interactive with Fiat--Shamir would additionally enter the random-oracle
model and require exact transcript binding and challenge sampling. This chapter
selects neither transformation yet.

Transparent succinct systems offer a different tradeoff. Spartan, for
example, provides argument-of-knowledge constructions for R1CS without a
relation-specific trusted setup in its transparent variants
\cite{setty2020spartan}. Its stack includes an arithmetic constraint system,
sum-check, and both computation and polynomial commitments. This is a credible
later production-oriented comparison point, but importing or reimplementing
such machinery would violate the present objective of a small, inspectable
next experiment.

Accordingly, MPC in the head is the leading family for a later explicitly
educational interactive experiment, conditional on circuit measurement and a
separate analysis of the full zero-knowledge and knowledge transformations.
It is not an implemented protocol, a selected production backend, or a current
security claim. Generic NP reductions and succinct arithmetic systems remain
comparison points; a trusted evaluator and private per-state computation remain
separate baselines.

\subsection{Completed Implementation Gate: A Fixed Game DAG}

The second executable evaluator still contains no proof system. It expresses
the public predicate as a fixed bottom-up directed acyclic graph and keeps the
recursive checker as an independent diagnostic baseline.

Independent finite enumeration gives 5,478 reachable states, of which 4,520
are nonterminal and 958 are terminal. The terminal states comprise 626 X wins,
316 O wins, and 16 draws. There are 8,631 legal edges out of X-turn states and
7,536 out of O-turn states, or 16,167 public state edges in total. By contrast,
the unmerged legal history tree has 549,946 nodes. The merged state graph is
therefore the natural fixed computation; these counts do not estimate a
cryptographic proof size. Independent tests also count 255,168 terminal legal
games in that unmerged tree.

The base-three board encoding supplies a particularly small topological
schedule. If cell $a$ is empty and the public side to move has numeric encoding
$t\in\{1,2\}$, then

\[
  \mathsf{child}(i,a)=i+t3^a>i.
\]

Consequently, one ascending scan of all 19,683 encodings can propagate
reachability and classify the graph, and one descending scan can evaluate
every child before its parent. A complete ascending scan between them validates
every policy byte. The implementation needs no recursion, queue, heap, edge
table, or stored ply schedule.

Canonical validation is part of the relation rather than a precondition hidden
from it. At every required state the evaluator compares the policy byte with
each public cell number. It never reads a board or successor at an address
chosen by that byte. It scans the full table, records the first malformed entry
only for diagnostics, continues the fixed evaluation, and gives the
invalid-policy verdict precedence over the root value. This prevents a missing
or nonsensical move from being misreported merely as a losing policy.

For every reachable state $i$, let $q_i$ denote whether the claimant avoids
loss from that state while following the policy. The recurrence is
well-founded from later plies to earlier plies; the implementation's
descending index scan is an equivalent reverse topological order because every
legal child index exceeds its parent index:

\[
q_i=
\begin{cases}
1, & i\text{ is terminal and is not an opponent win},\\
0, & i\text{ is an opponent win},\\
\displaystyle\bigvee_{a\in A(i)}
 \left([\pi(i)=a]\land q_{\mathsf{child}(i,a)}\right),
    & i\text{ is a claimant-turn state},\\
\displaystyle\bigwedge_{a\in A(i)}q_{\mathsf{child}(i,a)},
    & i\text{ is an opponent-turn state}.
\end{cases}
\]

Here $[\pi(i)=a]$ is a Boolean equality bit. Every state, legal move $a$, and
successor index is public; only the equality result and derived values depend
on the policy. Canonical-policy validation ensures that exactly one legal move
is selected at each claimant state. The $q_i$ values are deterministically
constrained intermediates, not independent semantic witness inputs; a backend
may nevertheless encode them as auxiliary wires or witness entries. The
recurrence accepts exactly when the empty-board value is one.

\begin{proposition}[DAG equivalence]
\label{prop:ttt-dag-equivalence}
For either claimant role, let $\pi$ be any Canonical Policy Version~1 table.
The fixed-DAG root value equals the Boolean verdict of the recursive checker.
\end{proposition}

\textbf{Proof.}
For a legal reachable state $s$, let $F(s)$ be the recursive semantic value
after canonical validation: it is zero exactly for an opponent-win terminal;
at a claimant turn it is the value of the single policy child; and at an
opponent turn it is the conjunction of every legal child value. We prove
$q(s)=F(s)$ by strong induction on the number of empty cells.

At a terminal state the two definitions agree directly. At a nonterminal
claimant state, canonical validity gives one legal move
$a^*=\pi(s)$. Every selector equality except $[\pi(s)=a^*]$ is zero, so the
disjunction reduces to $q(\mathsf{child}(s,a^*))$; the induction hypothesis
makes this $F(\mathsf{child}(s,a^*))=F(s)$. At a nonterminal opponent state,
both definitions are the conjunction of the values of all legal children, and
the induction hypothesis applies to each child. Therefore $q(s)=F(s)$ at every
reachable state and in particular at the empty board. $\square$

Malformed-table agreement is a separate validation statement. Both
implementations use the same required states, reason precedence, and first
ascending-index error. The differential tests compare that evidence, but do
not compare the recursive checker's first losing path or history-node counts
with a DAG report.

For claimant X, the C recurrence performs 8,631 policy-to-cell equalities,
8,631 selector term conjunctions, 8,631 selector disjunction folds, and 7,536
opponent conjunction folds. For claimant O, 7,536 replaces each selector
count and there are 8,631 opponent folds. Both roles scan 19,683 graph slots,
19,683 policy bytes, and 19,683 value slots; derive 5,478 reachable values; and
visit 16,167 public edges. These are operations of this C evaluator, not gates
in an optimized Boolean circuit.

The explicit scratch payload is two 19,683-byte arrays, one for state kinds and
one for values, plus a nine-byte board: 39,375 bytes in total. The serialized
policy remains exactly 19,683 bytes. This count excludes scalar locals, the
caller-owned report, compiler padding, and the call frame, so it is not a
whole-process memory measurement.

The differential corpus includes every legal Hamming-distance-one alternative
of the two generated reference policies: 6,208 for X and 5,439 for O. It also
includes 256 deterministically generated multi-entry policies per role and
malformed representatives. This does not enumerate the full canonical spaces,
whose sizes are

\[
  9\,7^{72}5^{756}3^{1372}
  \quad\text{for X},\qquad
  8^9 6^{252}4^{1140}2^{696}
  \quad\text{for O}.
\]

The multi-entry cases are a reproducible corpus, not a random sample, and no
probability bound is claimed. Tests can reveal implementation disagreement;
the induction gives the mathematical reason the two recurrences agree and may
still fail to expose a shared low-level C defect.

This equivalence gate supplies the semantics for the Boolean construction
below. The commitment computation remains separate rather than disappearing
inside the game predicate. Proof-internal view commitments would be an
additional cryptographic cost and assumption, not part of
\(\mathsf{Bind}\).

\subsection{Boolean Circuits for Core}

The implementation emits two circuits, \(\mathsf{Core\mbox{-}X}\) and
\(\mathsf{Core\mbox{-}O}\). Each specialization still receives and checks
the public role byte. The public input is the four-byte sequence

\begin{center}
\texttt{[version][role][initial-index-high][initial-index-low]}.
\end{center}

The index is an unsigned big-endian integer. Version~1 accepts only its
specialized role and initial index zero. The policy input is the unchanged
19,683-byte Canonical Policy Version~1 table. The IR calls that range
\emph{witness input}, but the evaluator receives every byte in plaintext.
Within each packed byte, bit zero is assigned first.

Wire zero is constant false and wire one is constant true. Public bits come
next, followed by policy bits and then gate outputs. Every gate has two earlier
wire operands and is either XOR or AND. Circuit shape, operands, and output wire
depend on the public game and role specialization, not on policy contents.
Required policy bytes reuse their legal-move equality wires between canonical
validation and the claimant recurrence; non-required bytes are compared with
\texttt{0xff}. The circuit follows the same public successor locations and
descending board-index recurrence as the fixed-DAG evaluator.

The canonical serialization is byte-defined rather than a dump of C structs.
It has a 32-byte header, big-endian integer fields, and one nine-byte record per
gate, for exactly \(32+9G\) bytes. Construction, evaluation, depth analysis,
and serialization use caller-owned storage. Evaluation uses one byte per wire
and applies a best-effort volatile wipe before returning. That wipe is not a
guarantee about registers, compiler-created copies, operating-system state, or
the caller-owned public and policy buffers.

\begin{table}[htbp]
\centering
\begin{tabular}{l@{\hspace{1.5em}}r@{\hspace{2em}}r}
\hline
Quantity & \(\mathsf{Core\mbox{-}X}\) & \(\mathsf{Core\mbox{-}O}\) \\
\hline
Public input bits & 32 & 32 \\
Policy/witness input bits & 157,464 & 157,464 \\
AND gates & 215,022 & 209,313 \\
XOR gates & 31,350 & 27,344 \\
Total gates & 246,372 & 236,657 \\
Total wires & 403,870 & 394,155 \\
Output depth & 31 & 30 \\
Gate storage (bytes) & 2,956,464 & 2,839,884 \\
Evaluator scratch (bytes) & 403,870 & 394,155 \\
Canonical serialization (bytes) & 2,217,380 & 2,129,945 \\
\hline
\end{tabular}
\caption{Exact Core circuit costs for the frozen Version~1 construction.}
\label{tab:ttt-core-cost}
\end{table}

These are exact counts for this frozen construction, confirmed by separately
analyzing and parsing the emitted circuits. They are not lower bounds,
optimized backend costs, proof sizes, transcript sizes, running-time results,
or security levels. The named builder scratch occupies 255,888 bytes on the
tested C11 ABI; this records \texttt{sizeof} for that layout rather than an
ABI-independent whole-process memory claim.

The bounded differential corpus covers the two reference policies, known
canonical losing policies, all 256 values at policy byte zero, named malformed
classes, 256 stratified legal one-entry mutations and 128 deterministic
canonical multi-entry policies per role, and invalid public version, role, and
initial-index inputs. Representative cases use the scalar production
evaluator. A separate test-only 64-lane interpreter reads the canonical bytes
independently so many policies can be compared with both preserved C
evaluators without changing the public API. This corpus is not exhaustive over
the enormous canonical policy spaces, and testing is not formal verification.

Serialization is not a commitment. Evaluating \(\mathsf{Core}\) does not
evaluate the complete relation
\(R_{\mathrm{V1}}=\mathsf{Core}\land\mathsf{Bind}\). The emitted Core circuits still contain
no commitment computation, proof-internal commitment, soundness, extraction,
zero knowledge, privacy, interaction, or production security.

\subsection{External Policy Commitment and Symbolic Bind Measurement}

Policy Commitment Profile Version~1 instantiates the external commitment
profile and disclosed-opening check in the future relation. It composes AC
Commitment Version~1 with the
fixed 16-byte protocol identifier
\texttt{AC-TTT-POLICY-V1}. Its commitment context contains protocol version
one, a public 32-byte session identifier, a public unsigned 32-bit statement
round, fixed protocol-role labels \texttt{PROVER} and \texttt{VERIFIER}, and
payload type one. The claimant's X/O game role is part of the payload and
therefore cannot be confused with the protocol-role labels.

The committed payload is exactly

\begin{center}
\texttt{[game version][claimant role][empty board u16be][policy:19683]}.
\end{center}

The public input proposed for the combined relation is 72 bytes:

\begin{center}
\texttt{[Core:4][session:32][statement round u32be][digest:32]}.
\end{center}

The witness-typed input is 19,715 bytes:

\begin{center}
\texttt{[canonical policy:19683][opening nonce:32]}.
\end{center}

The policy prefix is shared with \(\mathsf{Core}\). The canonical commitment
framing contributes 122 bytes around the 19,687-byte payload, so the hashed
message is 19,809 bytes. BLAKE2b uses 128-byte blocks and therefore performs
155 compression calls for this fixed length
\cite{saarinen2015blake2}. Protocol integers are encoded big-endian before
hashing; within the BLAKE2b computation, message words have the algorithm's
little-endian interpretation. The implementation uses the existing libsodium
generic-hash and random-byte interfaces
\cite{libsodium-generichash,libsodium-randombytes}.

The executable wrapper can create a fresh commitment, recompute one from a
caller-supplied nonce for test vectors, and verify a disclosed opening. The
published deterministic vector uses a synthetic, semantically invalid policy:
the external opening verifier accepts it, while the public policy checker
rejects it. This negative test is intentional. It satisfies
\(\mathsf{Bind}\)'s opening condition but not the complete relation.
Commitment consistency does not imply policy legality or non-loss, just as
\(\mathsf{Core}\) alone does not connect the accepted policy to \(C\).

The same security boundary as AC Commitment Version~1 applies. Correct
recomputation is tested. Hiding is only an informal random-oracle-style
argument requiring a fresh unpredictable secret 32-byte nonce. Binding is
only an informal collision/alternate-opening argument for the exact canonical
encoding. The profile provides no commit-time proof of knowledge,
authentication, identity, non-repudiation, fairness, delivery, availability,
or production guarantee. Session and round fields separate byte contexts only
if an application actually enforces fresh sessions and meaningful round
allocation.

Emitting a literal direct BLAKE2b Boolean circuit at this stage would create a
large artifact without yet supplying a private proof. The companion instead
walks the exact operation shape and wire depths of one frozen, deliberately
unoptimized XOR/AND representation. Each 64-bit addition uses

\[
\begin{aligned}
 c'&=((a\mathbin{\mathsf{XOR}}c)
       \mathbin{\mathsf{AND}}(b\mathbin{\mathsf{XOR}}c))
       \mathbin{\mathsf{XOR}}c,\\
 s&=(a\mathbin{\mathsf{XOR}}c)\mathbin{\mathsf{XOR}}b,
\end{aligned}
\]

Bit zero uses the one-AND/one-XOR half-adder implied by \(c=0\);
bits 1 through 62 use the relation above; and bit 63 emits only its two-XOR
sum because the final carry is discarded. This frozen construction therefore
uses 63 AND and 251 XOR gates per 64-bit addition. The model follows all
twelve rounds, rotations,
feed-forward operations, digest comparison, and the nonzero-session check.
It does not track bit values, emit a gate array, evaluate a circuit, or
serialize one. The deterministic commitment vectors separately check the
byte-level profile and the positive shared-policy path.

\begin{table}[htbp]
\centering
\begin{tabular}{lr}
\hline
Quantity & Symbolic \(\mathsf{Bind}\) model \\
\hline
Public input bits & 576 \\
Witness input bits & 157,720 \\
BLAKE2b compression calls & 155 \\
AND gates & 5,625,151 \\
XOR gates & 26,378,049 \\
Total gates & 32,003,200 \\
Wires, including two constants & 32,161,498 \\
Output depth & 1,558,999 \\
Flat 12-byte gate storage if emitted & 384,038,400 \\
One-byte-per-wire scratch if evaluated & 32,161,498 \\
Nine-byte-gate serialization if emitted & 288,028,832 \\
\hline
\end{tabular}
\caption{Exact symbolic Bind costs for the frozen direct representation.}
\label{tab:ttt-bind-cost}
\end{table}

Sharing the 576 public bits and the 157,720-bit witness layout, then adding one
final AND, gives the following purely arithmetical composition:

\begin{table}[htbp]
\centering
\begin{tabular}{lrr}
\hline
Quantity & \(\mathsf{Core\mbox{-}X}\land\mathsf{Bind}\)
         & \(\mathsf{Core\mbox{-}O}\land\mathsf{Bind}\) \\
\hline
AND gates & 5,840,174 & 5,834,465 \\
XOR gates & 26,409,399 & 26,405,393 \\
Total gates & 32,249,573 & 32,239,858 \\
Wires & 32,407,871 & 32,398,156 \\
Output depth & 1,559,000 & 1,559,000 \\
Serialization if emitted & 290,246,189 & 290,158,754 \\
\hline
\end{tabular}
\caption{Arithmetic composition of Core and Bind, without an emitted combined circuit.}
\label{tab:ttt-combined-cost}
\end{table}

These counts are exact only for the frozen operation shape. They are not lower
bounds, optimized-backend estimates, emitted circuits, equivalence proofs,
running-time measurements, proof sizes, transcript sizes, or security
parameters. The direct result is nevertheless informative:
\(\mathsf{Bind}\), not the finite game, dominates this Boolean route.
In this direct representation, Bind is about 130 to 135 times larger than Core
and contributes more than 99.2 percent of the combined gate count.
A proof-friendly commitment could change that tradeoff, but it would be a new
profile with new assumptions rather than an invisible optimization.

\subsection{Commitment Route Decision}

Policy Commitment Profile Version~1 remains the exact direct baseline. Its
API, vector, byte encoding, implementation, and symbolic counts remain
authoritative for disclosed-opening education and later comparison. The next
artifact will not materialize its 32,003,200-gate Bind circuit. This is a
sequencing decision, not a claim that BLAKE2b is insecure or that an unnamed
alternative is cheaper.

Nor is a replacement selected here. ``Proof-friendly'' is not the name of a
primitive or a security property. It is a measurable, backend-relative
requirement: a construction can be inexpensive in one field or constraint
language and expensive in another. Choosing it before the proof
representation would add assumptions and likely implementation machinery
without an end-to-end comparison. No new library is introduced by this
decision.

The next paper-first gate must compare source-backed pairs of proof
representations and commitments against the Version~1 baseline, then select
the two layers together. Any candidate pair must:

\begin{enumerate}
    \item preserve the intended temporal meaning: under its stated binding
    notion, one public value is intended to bind the same complete policy
    before the proof, across any declared repeated proofs, or at a later
    opening;
    \item define a new versioned profile, exact encoding, opening relation,
    setup, hiding and binding notions, and migration boundary;
    \item enforce the shared-witness condition above, rather than merely hash
    $C$ into a transcript or prove Core and Bind with independent policies;
    \item keep the external policy commitment distinct from any
    proof-internal commitment to randomized views or witness polynomials; and
    \item compare total relation, transcript, verifier, and proof costs under
    one stated security target, rather than compare only hash-gate counts.
\end{enumerate}

If no candidate survives that review, Version~1 remains the reference direct
route. If a candidate does survive, it must receive a new profile and relation
version; the existing 32-byte digest must never be silently reinterpreted.

No proof-internal view-commitment construction is selected. Conditioned only
on a future three-view shape using one 32-byte digest per view and one 32-byte
nonce per opened view, \(r\) repetitions would contribute \(96r\) digest
bytes and \(64r\) opened-nonce bytes. At the illustrative \(r=219\), this
conditional floor is 35,040 bytes. It excludes the opened views themselves,
challenges, metadata, transformations, the external policy commitment, and
every security reduction; it is not a proof-size estimate.
Because no proof-internal commitment profile is selected, this arithmetic is
recorded only here and is deliberately absent from the public C API.

\section{Repeated-Query Reconstruction}
\label{sec:ttt-repeated-query}

\begin{proposition}[Exact finite-oracle reconstruction]
\label{prop:ttt-oracle-reconstruction}
Let $D$ be a finite domain and let $\pi:D\rightarrow A$ be deterministic.
Any party permitted to query an oracle $O_\pi(s)=\pi(s)$ at every $s\in D$
can reconstruct $\pi$ using $|D|$ distinct queries.
\end{proposition}

\textbf{Proof.}
Enumerate the elements of $D$. Query each element once and store the returned
pair $(s,O_\pi(s))$. The resulting table agrees with $\pi$ at every member
of its domain and is therefore the same function. No computational assumption
is involved. $\square$

For Canonical Policy Version~1, unrestricted access over all claimant-turn
states reconstructs every required entry. If access is limited to games
actually reached from $s_0$, the proposition does not imply recovery of
unqueried off-path entries, even though those entries are legally reachable;
nevertheless, every observed move reveals the policy at one visited state.
Repeated adaptive games expose an expanding partial table.

A randomized policy changes exact recovery into statistical estimation:
repeated queries reveal samples from its output distribution. An anti-cloning
claim would then need a query budget, a target approximation metric, and a
definition of what counts as a successful clone.

One global zero-knowledge argument for the non-loss relation avoids serving a
per-state policy oracle. It can certify a property without offering the
strategy as a service. It also does not prevent the verifier from independently
constructing another non-losing Tic-Tac-Toe policy.

\section{Attacks, Abort, Leakage, and Composition}
\label{sec:ttt-attacks}

The public checker is designed to expose malformed and losing policies rather
than conceal them. Relevant cases include missing entries, out-of-range moves,
occupied-cell moves, stray bytes on opponent or terminal states, role
confusion, post-win boards, and a single losing line hidden behind many safe
branches.

Future private protocols add different risks:

\begin{itemize}
    \item a commitment may be replayed in a different game, role, or session
    if its context is incomplete;
    \item a malicious party may publish a syntactically valid value for which
    it knows no opening;
    \item a biased challenge or sampled checker may avoid the losing branch;
    \item selected openings or policy responses accumulate information across
    sessions;
    \item an interactive prover may abort before the verifier accepts;
    \item a recipient may learn a per-state output and then abort before
    providing its own promised action; and
    \item a protocol satisfying a stated zero-knowledge or knowledge-soundness
    property in a stand-alone execution may not retain that property under
    concurrent, reset, or adaptively interleaved executions.
\end{itemize}

No protocol can force a silent participant to send the next message. An abort
establishes non-completion, not motive, guilt, or network cause. Fairness,
guaranteed output delivery, availability, authentication, and recovery are
therefore not consequences of the future relation $R$. They would require
separate mechanisms and claims.

\section{Ethical and Institutional Interpretation}
\label{sec:ttt-ethics}

The attractive general idea is broader than this small game: a party may wish
to demonstrate that a private policy satisfies a public safety condition
without surrendering every private preference. The separation between policy
and proof can protect legitimate autonomy.

It can also hide power. A proof establishes only the encoded predicate. It
does not show that the predicate is fair, that excluded states were chosen
justly, that affected people can appeal, or that the hidden policy lacks other
harmful behaviour. In a consequential institution, the authority to define
the state space, allowed initial conditions, objective, proof circuit, and
update rules matters at least as much as witness privacy.

Tic-Tac-Toe is intentionally harmless and finite. It is a laboratory for this
distinction, not evidence that the same construction should govern people.

\section{What Cryptography Does Not Solve}
\label{sec:ttt-cryptography-does-not-solve}

Even a correct future zero-knowledge argument of knowledge for $R$ would not:

\begin{itemize}
    \item decide whether non-loss is the right objective;
    \item establish authorship, originality, identity, or exclusive knowledge;
    \item prove that the claimant will follow the policy after the proof;
    \item prevent independent discovery of another non-losing policy;
    \item hide outputs intentionally released by a policy oracle;
    \item force completion, explain an abort, or guarantee availability;
    \item make arbitrary losing initial conditions winnable; or
    \item give matrix inversion a role that has not been mathematically defined.
\end{itemize}

Zero knowledge can constrain what a transcript reveals about a witness. It
cannot choose a just predicate or turn a public, easily computed capability
into a unique strategic advantage.

\section{Limitations and Open Problems}
\label{sec:ttt-limitations}

The current result is limited to fixed ordinary three-by-three Tic-Tac-Toe, a
deterministic positional policy, one fixed role, and the standard empty board.
The complete policy is disclosed in memory. The external commitment wrapper
checks disclosed openings and has only the stated heuristic binding and hiding
arguments; verification receives the policy and nonce. The implementation contains no emitted
\(\mathsf{Bind}\) circuit, protocol-message parser, persistence, identity,
transport, proof system, private computation, random policy, or
generalized-game analysis.

The following questions remain open:

\begin{enumerate}
    \item Which public deterministic generation rule, if any, should connect
    the author's proposed session seed to a policy?
    \item Which initial positions may a verifier choose, and how should their
    pre-existing minimax values alter the claim?
    \item Which separately versioned commitment, if any, satisfies the
    required hiding and binding definitions while reducing the complete
    relation cost in the selected proof representation, and how does its total
    proof cost compare with the frozen Version~1 baseline?
    \item Which concrete MPC-in-the-head transformation can prove knowledge of
    the complete committed witness while meeting a malicious-verifier
    zero-knowledge definition; under which proof-internal commitment, setup,
    and composition assumptions; and with what rounds, transcript, and proof
    size?
    \item Should per-state evaluation reveal the query, the move, both, or only
    a derived game action?
    \item What query budget and leakage metric would make a meaningful
    anti-cloning claim?
    \item Does the author's matrix intuition correspond to a useful linear or
    algebraic representation, or should it remain a metaphor?
    \item Can a role-specific compact internal witness be expanded
    canonically to the complete Version~1 policy without changing the message
    bound by the commitment?
\end{enumerate}

These remain research questions. The recursive checker, fixed-DAG evaluator,
Boolean Core, and disclosed-opening wrapper are public baselines for later
comparisons.

\section{Exercises}
\label{sec:ttt-exercises}

\subsection*{Conceptual}
\begin{enumerate}
    \item For the board with X at cell~0 and O at cell~4, calculate the
    base-three board index and identify the policy-table entry used on X's
    turn.
    \item Identify the public input and witness in the existential and
    committed statement modes. Explain why a commitment is not a proof of
    legality or non-loss.
    \item Classify Core, Bind, an argument of knowledge, zero knowledge, and
    deployment using Table~\ref{tab:ttt-concept-ladder}.
    \item Give one replay risk and one abort risk that remain even if a future
    private argument for the relation is correct.
    \item Apply the concept ladder to the Poker chapter. State one place where
    the analogy holds and one place where the game introduces a new problem.
\end{enumerate}

\subsection*{Mathematical}
\begin{enumerate}
    \item Independently enumerate the legal reachable states and reproduce the
    four state counts in this chapter.
    \item Construct a canonical policy containing exactly one reachable losing
    branch, and explain why random play may miss it.
    \item Give an example showing why non-loss from $s_0$ does not imply
    locally optimal moves at every off-path state.
    \item Define a public policy-generation algorithm and extend $R$ so that
    a private seed, rather than a policy table, is the witness. List the new
    assumptions this introduces.
    \item Prove the finite-oracle reconstruction proposition for the role you
    choose, and state why it does not automatically reconstruct unqueried
    off-path states when access is restricted to actual games.
    \item Choose a legal state that is already losing for the player to move.
    Explain why unrestricted verifier-selected initial conditions make the
    original non-loss claim impossible.
    \item Describe the difference between a zero-knowledge argument for $R$
    and a zero-knowledge argument of knowledge for $R$. Identify what an
    extractor must output in the latter case.
    \item Construct two nonempty sets of policy witnesses whose intersection
    is empty. Use them to explain why separate Core and Bind arguments do not
    establish the committed-policy relation without explicit equality of the
    policy witnesses.
\end{enumerate}

\subsection*{Implementation}
\begin{enumerate}
    \item Prove directly from the base-three board encoding that every legal
    child index exceeds its parent index. Reimplement the bottom-up recurrence
    without calling either C evaluator and compare all three verdicts on a
    precisely described finite corpus.
    \item Independently parse, recount, and evaluate both serialized Core
    circuits. Explain why agreement on a bounded corpus proves neither policy
    privacy nor circuit equivalence for every possible policy.
\end{enumerate}

\subsection*{Research}
\begin{enumerate}
    \item Compare two candidate proof representations against the target
    experiments in Section~\ref{sec:ttt-security-target}. Account for setup,
    commitment compatibility, verifier work, transcript size, and composition.
    \item Propose an authenticated publication record for
    Table~\ref{tab:ttt-commitment-lifecycle}. State exactly which temporal and
    replay claims it supports and which it does not.
\end{enumerate}
